\section{外层问题与置换测试}
\label{sec:game}

现在用潜在 gadget 构造置换实例。其顶点编码奇偶方程元组上的仿射表。
当两张表在相同的保留坐标上诱导同一仿射映射（含同一截距）时，我们
把二者精确等同。这一公共表示将使我们在 Section~\ref{sec:matrix}
中比较方程元组与稀疏投影元组。

\subsection{带权奇偶方程与不同位置}

\emph{带权奇偶实例}由有限变量集 \(\mathcal X\)、有限的非空方程出现
集 \(\mathcal E\)、以及和为 1 的正有理权重 \(\omega_e\) 组成。一个
出现是一条有序方程
\[
 x_{v_{e,1}}+x_{v_{e,2}}+x_{v_{e,3}}=b_e,
 \qquad b_e\in\F_2.
\]
其值是被满足方程的加权比例，\(\OPT(\mathcal E)\) 记最大值。即使
有序方程重合，不同的出现也有不同的 ID。

\begin{theorem}[带权奇偶间隙]
\label{thm:parity}
对每个固定的有理 \(\xi>0\)，存在从布尔可满足性到带权奇偶实例的
确定性多项式时间归约，权重为二进制编码的有理数，使得
\[
 \begin{array}{ll}
 \textup{YES：}&\OPT(\mathcal E)\ge 1-\xi,\\
 \textup{NO：}&\OPT(\mathcal E)\le \frac12+\xi.
 \end{array}
\]
运行时间多项式可以依赖 \(\xi\)。
\end{theorem}

这是 \cite[Theorem~5.4 and its proof]{Hastad01} 中 H\aa stad 构造
的间隙推论。更具体地，该证明使用固定的二进验证者误差 \(\delta>0\)，
给出完备性 \(1-\delta\) 与可靠性 \((1+\delta)/2\)，并把验证者的
选择作为方程枚举出来、以其有理概率为权重。取 \(\delta\le\xi\) 即
得所示陈述。该枚举是确定性归约；验证者的随机选择描述了所得的带权
实例。

\begin{lemma}[使位置互异]
\label{lem:distinct-parity}
由带权奇偶实例 \(\mathcal E\) 可以在多项式时间内确定性构造实例
\(\mathcal E'\)：其中每个出现的三个变量 ID 互不相同。其权重是
多项式位长的正有理数，且
\[
 \OPT(\mathcal E')\ge\OPT(\mathcal E),
 \qquad
 \OPT(\mathcal E')
 \le \frac{1+\OPT(\mathcal E)}2+\frac3{100}.
\]
特别地，先取 \(\xi<1/100\) 的 Theorem~\ref{thm:parity}，再做该
预处理，给出完备性至少 \(1-\xi\)、NO 最优值至多 \(4/5\)。
\end{lemma}

\begin{proof}
把每个变量 \(v\) 换成一百个克隆 \((v,c)\)（\(c\in[100]\)）。对每
个原出现，均匀选取其三个克隆，只要各位置 underlying 变量相同就
强制选互异的克隆。列出所得的全部出现，并赋予相应的条件权重。每个
原出现至多有 \(100^3\) 种选择。给所有克隆其原变量的标记即证完备性
断言。

对上界，固定任意克隆标记，令
\[
 t_v=\frac1{100}\sum_{c=1}^{100}(-1)^{x_{(v,c)}}.
\]
先独立选克隆，不以互异性为条件。出现 \(e\) 的满足概率是
\[
 \frac{1+(-1)^{b_e}t_{v_{e,1}}t_{v_{e,2}}t_{v_{e,3}}}{2}.
\]
把每个偏差四舍五入为多数标记，零偏差任意处理。对被该取整赋值违背
的每条方程，分子中的带符号乘积非正，故其满足概率至多 \(1/2\)；
其余方程的概率至多 1。加权平均因此至多
\((1+\OPT(\mathcal E))/2\)。独立选取的克隆之间出现禁用碰撞的概率
由对三对位置的并界至多为 \(3/100\)。以其补为条件会把任何
\([0,1]\) 值函数的期望改变至多 \(3/100\)。界得证。除原出现权重外，
所有条件概率的分母都是常数规模。
\end{proof}

在构造的其余部分，\(\mathcal E\) 表示预处理后的实例，故每条方程
都有三个互异的变量 ID。其\emph{普通「方程对变量」博弈}按 \(\omega\)
采样 \(e\)、从 \([3]\) 均匀采样 \(j\)。第一个证明者收到出现 ID
\(e\)，以该方程的一个满足三元组作答；第二个收到变量 ID
\(v_{e,j}\)，以一个比特作答。接受表示位置 \(j\) 处相等。由于方程
内的 ID 互异，问题对决定了投影约束。

对任意第二证明者策略，其比特构成一个全局赋值。在该赋值违背的方程
上，每个合法的第一证明者答案至少在一个位置不一致。因此这个记作
\(G_{\mathcal E}\) 的普通投影博弈满足
\begin{equation}
 \val(G_{\mathcal E})
 \le 1-\frac{1-\OPT(\mathcal E)}3
 \le \frac{14}{15}
 \quad\text{若 }\OPT(\mathcal E)\le\frac45.
 \label{eq:base-game-gap}
\end{equation}
答案字母表规模分别为 4 与 2。

\subsection{元组、局部答案副本与齐次坐标}

固定正整数 \(k\)。方程问题 \(U=(e_1,\ldots,e_k)\) 从
\(\omega^{\otimes k}\) 有放回采样。其答案空间是 \(k\) 个仿射解平面
的乘积。块总是使用局部副本：即使同一方程或变量 ID 出现在多个块中，
它们的局部答案也是分离的坐标。

对参数 \(0<\beta<1\)，在每个块中独立地：以概率 \(1-\beta\) 保留该
方程，或以概率 \(\beta\) 只保留均匀所选位置上的变量 ID。所得问题
记为 \(O\)，\(\pi\) 是从 \(U\) 上的答案到 \(O\) 上的答案的坐标
投影。问题 \(O\) 在每个有序块位置记录一个方程 ID 或变量 ID。
单例块接受任一比特。

把方程答案用其前两个比特 \((x_i,y_i)\) 表示；第三个是
\(b_{e_i}+x_i+y_i\)。单例用其实际比特表示。加一个公共齐次坐标
\(t_0\)。于是
\[
 L_U=\F_2^{1+2k},\qquad m=1+2k,
\]
而 \(L_O\) 有一个齐次坐标、每个保留方程两个坐标、每个单例一个
坐标。在任一空间中，第一个坐标为 1 的向量恰是合法答案。把读取该
坐标的泛函记为 \(e_0^\top\)。

仿射投影 \(\pi\) 延拓为保持 \(t_0\) 的线性满射
\(\bar\pi:L_U\to L_O\)。在保留的方程块上它是恒等；在投影为单例的
块上，其坐标是
\begin{equation}
 \begin{cases}
 x_i,&\text{位置 }1,\\
 y_i,&\text{位置 }2,\\
 b_{e_i}t_0+x_i+y_i,&\text{位置 }3.
 \end{cases}
 \label{eq:outer-projection}
\end{equation}
把由这些逐块选择得到的每个映射称为\emph{容许投影}；这一类映射与
\(\beta\) 无关。满射性由固定 \(t_0\) 并在每个块中独立选取坐标
得到。对某些块被省略的问题，我们沿用同样的约定。

\subsection{精确表键与折叠}

跨投影问题等同表的做法遵循 Khot--Safra 的虚拟表方法
~\cite[Section~1.1]{KS13}。我们先把每张表用它所经过的坐标表示。
这将给投影表与其拉回相同的表示。

固定 Lemma~\ref{lem:latent} 的潜在数据 \((\mathcal V,K,C,\mu)\)，
其中 \(K=\F_2^\ell\)，并固定线性补 \(\mathcal V=K\oplus W\)。问题
\(Q\) 上的\emph{表}是从其答案空间到 \(\mathcal V\) 的仿射映射，
等价地其唯一线性延拓 \(P\in\Hom(L_Q,\mathcal V)\)。此处 \(Q\)
可以是方程问题或其某个容许投影，同一定义也适用于有省略块的问题。
把 \(P\) 中齐次坐标 \(t_0\) 的系数记为 \(c\)。

在方程块中，检查其斜率映射 \(\F_2^2\to\mathcal V\) 的秩，其贡献
记为 \(a_i x_i+d_i y_i\)。若秩为 2，保留该方程出现与该贡献。若秩
为 0，省略该块。剩余情形秩为 1，系数恰有
\[
 (a_i,d_i)=(v,0),\quad (0,v),\quad (v,v),
 \qquad v\ne0
\]
三种形式之一。这些形式分别只保留位置一、二、三处的变量，贡献为
实际比特 \(z_i\) 上的 \(v z_i\)。对第三种形式，
\[
 v(x_i+y_i)=v z_i+b_{e_i}v t_0,
 \qquad z_i=b_{e_i}t_0+x_i+y_i,
\]
故该替换同时给齐次截距加上 \(b_{e_i}v\)。在单例块中，恰在其斜率
非零时保留其变量 ID 与原贡献；否则省略该块。

保留的块构成 \((Q,P)\) 的\emph{支撑}。它们保持自己在 \([k]\) 中的
指标、方程/变量标签与 ID。保留的方程保持其有序的前两个坐标，保留
的单例用其实际比特坐标。于是不同块中的重复 ID 仍有分离的局部坐标。
支撑可以为空；无论何种情形齐次坐标都在。

设 \(J_3\) 是被约化到第三比特的方程块的集合，把这种块中的非零
斜率记为 \(v_i\)。支撑上的\emph{约化表} \(\widetilde P\) 具有上述
保留块贡献，及齐次截距
\begin{equation}
 \widetilde c=c+\sum_{i\in J_3} b_{e_i}v_i.
 \label{eq:reduced-table-intercept}
\end{equation}
为核对完整表，设 \(L_{\rm supp}\) 为支撑的齐次答案空间，
\(\rho_{Q,P}:L_Q\to L_{\rm supp}\) 保留 \(t_0\) 与指定的块坐标。
它在保留的方程上、以及在 \(Q\) 中已有的保留单例块上是恒等；在
被约化的方程上使用 Equation~\eqref{eq:outer-projection} 的相应
实际比特坐标；并丢弃被省略的块。在每个块中独立选坐标可知
\(\rho_{Q,P}\) 满射。各块恒等式与
Equation~\eqref{eq:reduced-table-intercept} 给出
\[
 P=\widetilde P\rho_{Q,P}.
\]
满射性使 \(\widetilde P\) 成为具有此分解的唯一线性映射。我们把
\((Q,P)\) 的\emph{键}定义为其支撑连同这整张约化表。于是两张表有
相同键，当且仅当它们的保留坐标一致（相同的位置、标签与 ID），且
其诱导仿射映射在这些坐标上一致（包括截距）。

接下来核对与每个容许投影 \(\bar\pi:L_U\to L_O\) 的相容性。设
\(P'\in\Hom(L_O,\mathcal V)\)。\(O\) 中保留的方程在 \(P'\) 及其
拉回 \(P'\bar\pi\) 中有相同的贡献，故其约化相同。斜率为零的单例在
两侧都消失。斜率非零 \(v\) 的单例拉回为上述三种秩一模式之一，约化
到同一块位置上的同一变量 ID。对第三位置的单例，拉回给截距加上
\(b_{e_i}v\)，而其约化再加一次 \(b_{e_i}v\)。两项在 \(\F_2\) 上
相消。其余截距修正都来自保留的方程块，且在两侧一致。因此支撑与
完整约化表一致：\((U,P'\bar\pi)\) 与 \((O,P')\) 有相同的键。

把键的全域定义为方程问题 \(U\in\mathcal E^k\) 上所有表的键。刚证
的相容性表明：每个容许投影问题上的表的键都属于这个全域——从任何
相容的 \(U\) 拉回即可。

现在应用 Bellare、Goldreich 与 Sudan 的折叠原理
~\cite[Section~3.3]{BGS98}：一个代表元上的取值决定其平移上的取值。
给键加常数 \(h\in K\) 指给其约化截距加 \(h\)，支撑与斜率不动。
这由把呈现表 \(P\) 换成 \(P+h e_0^\top\) 实现，故保持键全域。
由于键记录了截距，该作用是自由的。我们的置换实例的顶点就是该作用
的一条轨道，由其唯一的、约化截距 \(K\) 分量为零的键表示。

对方程问题或容许投影问题 \(Q\) 上的表 \(P\)，令 \(c(Q,P)\in K\)
为其键的约化截距的 \(K\) 分量，\(v(Q,P)\) 为去掉该分量后的键所
代表的轨道顶点。顶点的染色 \(\lambda\) 定义恢复表答案
\[
 F_Q(P)=\lambda(v(Q,P))+c(Q,P).
\]
加常数只改变被去掉的分量，故
\begin{equation}
 F_Q(P+h e_0^\top)=F_Q(P)+h
 \qquad(h\in K).
 \label{eq:table-folding}
\end{equation}
同键的表有相同的代表与被去分量。因此拉回相容性给出：对每个容许的
\(U,O,\bar\pi\) 与 \(O\) 上的每张表 \(P'\)，
\begin{equation}
 F_U(P'\bar\pi)=F_O(P').
 \label{eq:key-sharing}
\end{equation}
于是投影问题使用由同一顶点染色决定的答案，无需额外的一致性测试。

\subsection{潜在矩阵测试与完备性}

置换实例记作 \(\mathcal G\)，独立采样
\[
 U\sim\omega^{\otimes k},\qquad
 P\text{ 在 }\Hom(L_U,\mathcal V)\text{ 中均匀},\qquad
 a\sim\mu,\qquad l\text{ 在 }\F_2^m\text{ 中均匀},
\]
并检验
\begin{equation}
 F_U(P)=F_U(P+a l^\top).
 \label{eq:matrix-test}
\end{equation}
此处 \(a l^\top\) 表示映射 \(z\mapsto a(l^\top z)\)。每个样本指定
两个键轨道顶点之间的一条边。若其恢复偏移为 \(c_0,c_1\)，则其约束
是固定字母表 \(K\) 上的置换
\[
 \lambda(v_1)=\lambda(v_0)+c_0+c_1.
\]
边出现以采样概率为权重。此阶段允许自环与重复边。涉及 \(O\) 的
稀疏投影实验用于分析这些同样的顶点标记；它不是 \(\mathcal G\) 中
的额外测试。

\begin{lemma}[外层矩阵测试的完备性]
\label{lem:outer-completeness}
若 \(\mathcal E\) 有值至少 \(1-\xi\) 的赋值，且潜在映射满足
\[
 \Pr_{x\text{ 在 }\mathcal V\text{ 中均匀},\ a\sim\mu}
       [C(x+a)\ne C(x)]\le p_*,
\]
则
\[
 \val(\mathcal G)\ge 1-k\xi-\frac{p_*}{2}.
\]
\end{lemma}

\begin{proof}
固定近满足的全局比特赋值。给每个约化支撑如下合法答案：单例用其
全局比特；保留的方程用其前两个全局比特，并置第三个为满足该方程的
唯一值。这些选择只依赖支撑（包括其出现 ID），因此对每个精确键
独立于其来源地定义了一个答案。

在轨道代表元处，用该答案对其约化表求值并应用 \(C\)。这定义一个
顶点染色。\(C\) 在 \(K\) 下的等变性蕴含：其恢复答案是对实际约化表
求值再作用 \(C\)，与 Equation~\eqref{eq:table-folding} 一致。

以至少 \(1-k\xi\) 的概率，\(U\) 中的每条方程都被全局赋值满足。
在这样的元组上，\emph{每张}表使用的支撑答案都是同一个完整答案
\(z_U\in L_U\)（首坐标为 1）的投影。于是
Equation~\eqref{eq:matrix-test} 的两侧是
\[
 C(Pz_U)\quad\text{与}\quad
 C(Pz_U+a(l^\top z_U)).
\]
由于 \(z_U\ne0\)，向量 \(Pz_U\) 在 \(\mathcal V\) 中均匀，与
\(a,l\) 独立，而 \(l^\top z_U\) 是独立的公平比特。因此测试在这些
元组上以至多 \(p_*/2\) 的概率失败。加上元组含未满足方程的概率即
得该界。
\end{proof}
